\documentclass[12pt,a4paper,oneside]{book}
\usepackage[utf8]{inputenc}
\usepackage[T1]{fontenc}
\usepackage[english,french]{babel}
\usepackage{graphicx}
\usepackage{geometry}
\usepackage[table]{xcolor}
\usepackage{colortbl}
\usepackage{titlesec}
\usepackage{fancyhdr}
\usepackage{hyperref}
\usepackage{xcolor}
\usepackage{listings}
\usepackage{multicol}
\usepackage{amsmath}
\usepackage{amssymb}
\usepackage{booktabs}
\usepackage{longtable}
\usepackage{parskip}
\usepackage{enumitem}
\usepackage{fontawesome5}
\usepackage{url}

\definecolor{electronblue}{RGB}{0, 120, 212}
\definecolor{electronpurple}{RGB}{136, 23, 152}
\definecolor{electroncyan}{RGB}{0, 190, 255}
\definecolor{electrongold}{RGB}{255, 185, 0}
\definecolor{darkbg}{RGB}{15, 23, 42}
\definecolor{lightbg}{RGB}{249, 250, 251}
\definecolor{softgray}{RGB}{248, 249, 250}
\definecolor{mediumgray}{RGB}{107, 114, 128}
\definecolor{success}{RGB}{34, 197, 94}
\definecolor{warning}{RGB}{251, 191, 36}

\geometry{top=2.5cm,bottom=2.5cm,left=2.5cm,right=2.5cm}

\hypersetup{
    colorlinks=true,
    linkcolor=electronblue,
    filecolor=electronpurple,
    urlcolor=electroncyan,
    pdfauthor={E-Graphisme by ELECTRON},
    pdftitle={Documentation Enterprise - Volume 5},
    pdfsubject={AI Center & Agents IA Enterprise}
}

\pagestyle{fancy}
\fancyhf{}
\fancyhead[LE,RO]{\thepage}
\fancyhead[RE,LO]{\leftmark}
\fancyfoot[CE,CO]{\thepage}

\titleformat{\chapter}[display]
{\normalfont\Huge\bfseries\color{electronblue}}
{\chaptertitlename\ \thechapter}{20pt}{\Huge}
\titleformat{\section}
{\normalfont\Large\bfseries\color{electronblue}}
{\thesection}{1em}{}
\titleformat{\subsection}
{\normalfont\large\bfseries\color{electronpurple}}
{\thesubsection}{1em}{}

\begin{document}

\begin{titlepage}
\thispagestyle{empty}
\begin{center}
\vspace*{2cm}
{\Huge\bfseries\color{electronblue}E-GRAPHISME}\\[0.3cm]
{\LARGE\bfseries\color{electronpurple}by ELECTRON}\\[2cm]
\vspace{2cm}
{\huge\bfseries\color{darkbg}AI CENTER}\\[0.3cm]
{\huge\bfseries\color{darkbg}Le Centre de Commandement de l'Intelligence Artificielle}\\[1cm]
{\Large\bfseries\color{electroncyan}Volume 5}\\[3cm]
\vspace{2cm}
{\Large\color{mediumgray}Une entreprise où chaque département est assisté par une IA spécialisée,}\\[0.5cm]
{\Large\color{mediumgray}coordonnée depuis un centre de supervision privé réservé au Fondateur.}\\[3cm]
{\Large\bfseries\color{electrongold}Version 2026}\\[2cm]
\vfill
\end{center}
\end{titlepage}

\frontmatter

\chapter*{Préface}
\addcontentsline{toc}{chapter}{Préface}
Ce volume présente le AI Center et les Agents IA Enterprise, le centre de commandement de l'intelligence artificielle de la plateforme E-Graphisme by ELECTRON.

\bigskip
\textbf{Classification :} Document Interne - Partageable\\
\textbf{Volume :} 5\\
\textbf{Version :} 2026

\tableofcontents

\mainmatter

\chapter{LE AI CENTER}

\section{Présentation}
Le AI Center constitue le cerveau numérique de E-Graphisme by ELECTRON. Tous les agents IA résident dans cet environnement sécurisé.

Contrairement à un chatbot classique, il s'agit d'une organisation virtuelle, où chaque IA possède un métier, des responsabilités, une mémoire, des outils et un domaine de compétence.

Toutes les interactions sont exécutées localement grâce à Ollama, garantissant confidentialité et maîtrise des données.

\chapter{ARCHITECTURE GÉNÉRALE}

\section{Fonctions de l'Orchestrator}
Le centre est organisé autour d'un AI Orchestrator, qui coordonne les échanges entre les agents.

\bigskip
\begin{center}
\begin{tabular}{|c|p{10cm}|}
\hline
\rowcolor{electronblue}
\color{white}\textbf{Fonction} & \color{white}\textbf{Description} \\
\hline
Répartition des tâches & Distribution aux agents adaptés \\
\hline
Planification & Ordonnancement des actions \\
\hline
Priorisation & Gestion des urgences \\
\hline
Suivi d'exécution & Monitoring des tâches \\
\hline
Journalisation & Logging complet \\
\hline
Gestion des événements & Réponse aux incidents \\
\hline
Synchronisation des mémoires & Partage d'informations \\
\hline
Remontée des alertes & Notification des problèmes \\
\hline
\end{tabular}
\end{center}

\chapter{AI ORCHESTRATOR}

\section{Rôle}
L'AI Orchestrator remplace les moteurs d'automatisation externes. Il reçoit les demandes, identifie les compétences nécessaires, sollicite les agents concernés puis consolide leurs réponses.

\section{Exemple de Flux}
\begin{enumerate}
\item Un client demande un site e-commerce
\item L'Orchestrator consulte le Commercial IA
\item Le Designer IA prépare une proposition graphique
\item Le Développeur IA évalue la faisabilité
\item Le Finance IA établit un devis
\item Le CEO AI valide si nécessaire
\item Les tâches sont créées dans l'ERP
\end{enumerate}

\chapter{CENTRE PRIVÉ DU FONDATEUR}

\section{Accès}
Accessible uniquement au Fondateur.

\section{Fonctionnalités}
Cette interface permet de :
\begin{itemize}
\item voir tous les agents ;
\item démarrer ou arrêter un agent ;
\item choisir le modèle Ollama utilisé par chaque agent ;
\item modifier les prompts système ;
\item consulter les journaux ;
\item gérer les connaissances ;
\item surveiller les performances CPU/GPU ;
\item suivre les conversations internes.
\end{itemize}

\chapter{HIÉRARCHIE DES AGENTS}

\section{Niveau Exécutif}
\begin{center}
\begin{tabular}{|c|p{10cm}|}
\hline
\rowcolor{electronpurple}
\color{white}\textbf{Agent} & \color{white}\textbf{Rôle} \\
\hline
CEO AI & Direction générale \\
\hline
Executive Assistant AI & Assistant exécutif \\
\hline
\end{tabular}
\end{center}

\section{Niveau Direction}
\begin{multicols}{2}
\begin{itemize}
\item Commercial AI
\item Marketing AI
\item Finance AI
\item RH AI
\item Juridique AI
\item Production AI
\end{itemize}
\end{multicols}

\section{Niveau Créatif}
\begin{multicols}{2}
\begin{itemize}
\item Designer AI
\item Branding AI
\item Art Director AI
\item Motion AI
\item Interior Designer AI
\end{itemize}
\end{multicols}

\section{Niveau Technique}
\begin{multicols}{2}
\begin{itemize}
\item Frontend AI
\item Backend AI
\item Full Stack AI
\item DevOps AI
\item Database AI
\item Cybersecurity AI
\end{itemize}
\end{multicols}

\section{Niveau Opérationnel}
\begin{multicols}{2}
\begin{itemize}
\item Support AI
\item Marketplace AI
\item Impression AI
\item Logistique AI
\item Community Manager AI
\item Formation AI
\end{itemize}
\end{multicols}

\chapter{DÉPARTEMENTS IA}

\section{CEO AI}
\subsection{Responsabilités}
\begin{itemize}
\item supervision globale ;
\item arbitrage ;
\item validation des décisions importantes ;
\item suivi des objectifs.
\end{itemize}

\section{Executive Assistant AI}
\subsection{Fonctions}
\begin{itemize}
\item planification ;
\item agenda ;
\item rappels ;
\item préparation des réunions ;
\item synthèse des rapports.
\end{itemize}

\section{Commercial AI}
\subsection{Fonctions}
\begin{itemize}
\item qualification des prospects ;
\item devis ;
\item suivi commercial ;
\item relances ;
\item gestion des contrats.
\end{itemize}

\section{Marketing AI}
\subsection{Fonctions}
\begin{itemize}
\item stratégie digitale ;
\item campagnes ;
\item SEO ;
\item réseaux sociaux ;
\item génération de contenus.
\end{itemize}

\section{Designer AI}
\subsection{Fonctions}
\begin{itemize}
\item création de logos ;
\item affiches ;
\item identités visuelles ;
\item supports imprimés ;
\item maquettes.
\end{itemize}

\section{Art Director AI}
\subsection{Fonctions}
\begin{itemize}
\item supervision artistique ;
\item contrôle qualité ;
\item cohérence visuelle ;
\item validation des créations.
\end{itemize}

\section{Interior Designer AI}
\subsection{Fonctions}
\begin{itemize}
\item recommandations de décoration ;
\item simulation d'intégration d'œuvres ;
\item harmonie des couleurs ;
\item choix des cadres.
\end{itemize}

\section{Frontend AI}
\subsection{Fonctions}
\begin{itemize}
\item interfaces React ;
\item composants ;
\item animations ;
\item responsive design.
\end{itemize}

\section{Backend AI}
\subsection{Fonctions}
\begin{itemize}
\item API ;
\item authentification ;
\item logique métier ;
\item sécurité.
\end{itemize}

\section{Full Stack AI}
\subsection{Fonctions}
\begin{itemize}
\item coordination entre Frontend et Backend ;
\item optimisation des performances.
\end{itemize}

\section{DevOps AI}
\subsection{Fonctions}
\begin{itemize}
\item Docker ;
\item sauvegardes ;
\item surveillance ;
\item déploiements.
\end{itemize}

\section{Database AI}
\subsection{Fonctions}
\begin{itemize}
\item PostgreSQL ;
\item Redis ;
\item Firebase ;
\item optimisation des requêtes.
\end{itemize}

\section{Cybersecurity AI}
\subsection{Fonctions}
\begin{itemize}
\item contrôle des accès ;
\item audit ;
\item chiffrement ;
\item surveillance.
\end{itemize}

\section{Finance AI}
\subsection{Fonctions}
\begin{itemize}
\item devis ;
\item factures ;
\item paiements ;
\item rapports financiers.
\end{itemize}

\section{RH AI}
\subsection{Fonctions}
\begin{itemize}
\item recrutement ;
\item gestion des collaborateurs ;
\item formation ;
\item évaluations.
\end{itemize}

\section{Support AI}
\subsection{Fonctions}
\begin{itemize}
\item assistance ;
\item tickets ;
\item FAQ ;
\item résolution des incidents.
\end{itemize}

\section{Marketplace AI}
\subsection{Fonctions}
\begin{itemize}
\item gestion des œuvres ;
\item publication ;
\item modération ;
\item recommandations.
\end{itemize}

\section{Impression AI}
\subsection{Fonctions}
\begin{itemize}
\item préparation des fichiers ;
\item contrôle des formats ;
\item gestion des commandes d'impression.
\end{itemize}

\section{Community Manager AI}
\subsection{Fonctions}
\begin{itemize}
\item calendrier éditorial ;
\item publications ;
\item interactions avec la communauté.
\end{itemize}

\chapter{COMMUNICATION ENTRE LES AGENTS}

\section{Messagerie Interne}
Les agents échangent des informations via une messagerie interne.

\section{Characteristics}
Chaque message est :
\begin{itemize}
\item horodaté ;
\item signé ;
\item enregistré ;
\item traçable.
\end{itemize}

Le Fondateur peut consulter ces échanges pour comprendre les décisions prises.

\chapter{MÉMOIRE DES AGENTS}

\section{Niveaux de Mémoire}
Chaque agent dispose de plusieurs niveaux de mémoire :
\begin{center}
\begin{tabular}{|c|p{10cm}|}
\hline
\rowcolor{electroncyan}
\color{white}\textbf{Type} & \color{white}\textbf{Description} \\
\hline
Mémoire de session & Données temporaires \\
\hline
Mémoire métier & Compétences et procédures \\
\hline
Mémoire documentaire & Base de connaissances \\
\hline
Historique des interactions & Traçabilité \\
\hline
Bibliothèque de procédures & Méthodologies \\
\hline
\end{tabular}
\end{center}

Les connaissances sont isolées par domaine afin d'éviter les interférences.

\chapter{BIBLIOTHÈQUE DE CONNAISSANCES}

\section{Bases Documentaires}
Chaque département possède une base documentaire dédiée.

\section{Exemples}
\begin{itemize}
\item normes graphiques ;
\item guides de développement ;
\item procédures commerciales ;
\item modèles de devis ;
\item chartes internes ;
\item FAQ.
\end{itemize}

\chapter{GESTION DES MODÈLES OLLAMA}

\section{Options du Fondateur}
Le Fondateur peut :
\begin{itemize}
\item associer un modèle à chaque agent ;
\item définir des profils de performance (rapide, équilibré, qualité) ;
\item mettre à jour les modèles ;
\item suivre leur utilisation et leur consommation de ressources.
\end{itemize}

\chapter{DASHBOARD DU FONDATEUR}

\section{Vue d'Ensemble en Temps Réel}
\begin{multicols}{2}
\begin{itemize}
\item état des agents ;
\item charge CPU/GPU ;
\item projets actifs ;
\item revenus ;
\item ventes Marketplace ;
\item statistiques ERP ;
\item alertes ;
\item sauvegardes ;
\item notifications.
\end{itemize}
\end{multicols}

\chapter{DASHBOARD DES AGENTS}

\section{Espace de Chaque Agent}
Chaque agent possède un espace avec :
\begin{multicols}{2}
\begin{itemize}
\item tâches en cours ;
\item historique ;
\item indicateurs de performance ;
\item outils disponibles ;
\item paramètres.
\end{itemize}
\end{multicols}

\chapter{JOURNAL DES ACTIVITÉS}

\section{Opérations Enregistrées}
Toutes les opérations importantes sont enregistrées :
\begin{multicols}{2}
\begin{itemize}
\item connexions ;
\item créations ;
\item modifications ;
\item validations ;
\item erreurs ;
\item actions des agents.
\end{itemize}
\end{multicols}

Ces journaux facilitent l'audit et le diagnostic.

\chapter{SÉCURITÉ ET PERMISSIONS}

\section{Contrôle d'Accès}
Le Centre IA applique un contrôle d'accès par rôles.

\section{Profils Principaux}
\begin{center}
\begin{tabular}{|c|p{10cm}|}
\hline
\rowcolor{electronblue}
\color{white}\textbf{Profil} & \color{white}\textbf{Accès} \\
\hline
Fondateur & Complet \\
\hline
Administrateur & Étendu \\
\hline
Employé & Restreint \\
\hline
Client & Limité \\
\hline
Affilié & Minimal \\
\hline
\end{tabular}
\end{center}

Chaque profil dispose uniquement des autorisations nécessaires.

\chapter{ÉVOLUTIONS FUTURES}

\section{Fonctionnalités Prévues}
Les prochaines versions pourront intégrer :
\begin{itemize}
\item assistants vocaux multilingues ;
\item réunions IA collaboratives ;
\item génération de vidéos et de présentations ;
\item analyse prédictive avancée ;
\item jumeau numérique de l'entreprise pour simuler des scénarios de croissance.
\end{itemize}

\appendix

\chapter{Index}
\begin{enumerate}
\item Présentation du AI Center
\item Architecture générale
\item AI Orchestrator
\item Centre privé du Fondateur
\item Hiérarchie des Agents IA
\item Départements IA
\item Communication entre les Agents
\item Mémoire des Agents
\item Bibliothèque de connaissances
\item Gestion des modèles Ollama
\item Dashboard du Fondateur
\item Dashboard des Agents
\item Journal des activités
\item Sécurité et permissions
\item Évolutions futures
\end{enumerate}

\end{document}
